\section{Test}
\label{sec:test}

\subsection{Teststrategie}
Die Teststrategie umfasste:

\begin{itemize}
    \item \textbf{Unit-Tests}: Für die Kernkomponenten (z. B. Statistik-Berechnung, Config-Handling, Measurer-Interface)
    \item \textbf{Integrationstests}: Für die Datenbankanbindung und API-Endpunkte
    \item \textbf{E2E-Tests}: Für das Web-Interface und PDF-Export
    \item \textbf{Manuelle Tests}: Für UI-Interaktionen und Benutzerfreundlichkeit
\end{itemize}

\subsection{JUnit-Testabdeckung}
SignalReport wurde mit JUnit 5 umfassend getestet. Die 18 Testklassen mit laufenden Tests -- in Paketen, die die Quellstruktur spiegeln (\texttt{config}, \texttt{measurement}, \texttt{network}, \texttt{storage}, \texttt{report}, \texttt{web}, \texttt{i18n}) -- mit insgesamt 156 Tests (zuzüglich zweier netz- bzw.\ PDF-gebundener Smoke-Klassen, die per Umgebungsvariable aktivierbar sind) decken folgende Bereiche ab:

\begin{itemize}
    \item \textbf{MeasurementTest} (5 Tests): Grundlegende Objekt-Erstellung und Validierung
    \item \textbf{HostIdentifierTest} (4 Tests): Stabile Host-Identifikation über Neustarts
    \item \textbf{StatisticsTest} (8 Tests): Mathematisch korrekte Berechnung von 95. Perzentil, Jitter und Paketverlust mit realen DB-Abfragen
    \item \textbf{ConfigTest} (17 Tests): Laden/Speichern von Konfiguration, hashPassword, Auth-Verifikation, Theme, Save+Reload
    \item \textbf{H2MeasurementRepositoryTest} (10 Tests): Vollstaendige Datenbank-Integration mit IP-Tracking, Host-Registrierung
    \item \textbf{MeasurerInterfaceTest} (6 Tests): Interface-Implementierung, Polymorphie, reale Messungen (Ping/DNS/HTTP)
    \item \textbf{MaintenanceWindowTest} (7 Tests): Deaktiviert, ganztägig, aktuelle Zeit, Mitternachtsübergang, Standardwerte
    \item \textbf{SessionManagerTest} (19 Tests): Nonce-Generierung und -Validierung (Single-Use, Ablauf), Session-Erstellung und -Timeout (24\,h), Rollen\-unter\-schei\-dung (Admin/User), Challenge-Response-Verifikation, Replay-Schutz
    \item \textbf{I18nTest} (10 Tests): Schlüssel-Parität aller neun Sprachdateien gegen die Referenz (de.json), Existenz aller im Quellcode verwendeten Platzhalter (inkl. \texttt{app.js}), Fallback-Verhalten, Sprach-Erkennung
    \item \textbf{GatewayDiscoveryTest} (15 Tests): Traceroute-Parsing, RFC-1918-Klassifizierung, Gateway-Kette (nah/fern), Erkennung virtueller Gateways (Docker-Bridge, VM-NAT) und Überspringen der Docker-Bridge
    \item \textbf{ConnectivityAssessmentTest} (8 Tests): Schuld-Verdikt (Router / Internet-Gateway / Internet) auf Basis des Paketverlusts je Segment
    \item \textbf{ReliabilityReportTest} (10 Tests): lückenbewusste Verfügbarkeit/Abdeckung/MTBF/MTTR, Aggregation aufeinanderfolgender Fehlmessungen zu einem Ausfall, Ausschluss via \texttt{excluded}-Spalte
    \item \textbf{ServiceReachabilityConfigTest} (8 Tests): Defaults (deaktiviert, 6-Stunden-Takt), Intervall-Grenzen, Pro-Dienst-Schalter, JSON-Round-Trip
    \item \textbf{ServiceReachabilityAssessmentTest} (15 Tests): Verdikt-Klassifizierung (erreichbar / Dienst-Störung / DNS-, TCP-, SNI-Sperre / Sperrseite) aus synthetischen Beobachtungen
    \item \textbf{ServiceReachabilityReportTest} (4 Tests): Verdichtung der Prüfungen zu Episoden (die Sperr-/Störungs-Timeline)
    \item \textbf{ServiceReachabilitySchedulerTest} (5 Tests): Leitungs-Gate-Entscheidung (nur erfolgreiche PING-/HTTP-Messungen zählen als Internet-Beleg)
    \item \textbf{ServiceCheckRepositoryTest} (3 Tests): \texttt{service\_checks}-Persistenz, Zeitraum-Abfrage, neueste Prüfung je Dienst
    \item \textbf{ServiceReachabilityRoutesTest} (2 Tests): \enquote{Jetzt prüfen}-Endpoint inkl.\ Abkühlphase (gegen echten Javalin-Server)
\end{itemize}

Alle Tests laufen ohne externe Abhängigkeiten (H2 In-Memory-Datenbanken für Integrationstests) und sind in unter 2 Sekunden durchführbar.

\subsection{Beispiel-Test für Statistik-Berechnung}
\begin{lstlisting}[language=Java, caption=Test fuer 95. Perzentil]
@Test
void testP95CalculationWithSortedLatencies() {
    // Arrange: 20 Messungen (95. Perzentil = 19. Wert)
    List<Double> latencies = new ArrayList<>();
    for (int i = 1; i <= 20; i++) {
        latencies.add((double) i * 10); // 10, 20, 30, ..., 200
    }

    // Act: 95. Perzentil = Index 18 (0-basiert) = 190.0
    int p95Index = (int) Math.ceil(latencies.size() * 0.95) - 1;
    double p95Value = latencies.get(p95Index);

    // Assert
    assertEquals(18, p95Index, "Index fuer 95. Perzentil muss 18 sein");
    assertEquals(190.0, p95Value, 0.01, "95. Perzentil-Wert muss 190.0 sein");
}
\end{lstlisting}

\subsection{Testergebnisse}
\begin{table}[H]
\centering
\small
\begin{tabular}{|p{8.5cm}|l|l|}
\hline
\textbf{Testkategorie} & \textbf{Anzahl} & \textbf{Erfolgsquote} \\
\hline
Unit-Tests (Measurement, HostIdentifier, Config, Measurer, MaintenanceWindow, SessionManager, I18n, GatewayDiscovery, ConnectivityAssessment, ReliabilityReport, ServiceReachabilityConfig, ServiceReachabilityAssessment, ServiceReachabilityReport, ServiceReachabilityScheduler) & 133 & 100\% \\
\hline
Integrationstests (Statistics, H2MeasurementRepository, ServiceCheckRepository, ServiceReachabilityRoutes) & 23 & 100\% \\
\hline
E2E-Tests (manuell: Web-UI, PDF-Export, CSV-Export) & 5 & 100\% \\
\hline
Smoke-Tests (opt-in: Netz-Probe, PDF-Erzeugung) & 3 & -- \\
\hline
\textbf{Gesamt} & \textbf{156 + 5} & \textbf{100\%} \\
\hline
\end{tabular}
\caption{Testergebnisse}
\label{tab:test}
\end{table}

\subsection{Testumgebung}
\begin{itemize}
    \item \textbf{Betriebssysteme}: Windows 10 64-Bit, Windows 11, Windows 11 on ARM,
      Raspberry OS, Debian 13.4, Ubuntu 25.10, macOS Tahoe 26, Synology DSM 7.3.2,
      Docker Desktop 4.67.0
    \item \textbf{Java-Version}: 21
    \item \textbf{Browser}: Chrome, Chromium, Vivaldi, Firefox, Edge, Opera, Safari
    \item \textbf{Hardware}: Desktop-PC (Intel Core i7 14.~Gen.), Desktop-PC
      (AMD Ryzen 5 7040 Series), Notebook (Intel Core Ultra i7 100~Serie),
      Raspberry Pi 4, Raspberry Pi 5, Synology NAS DS420j, MacBook Air M2
\end{itemize}

\subsection{Test-Philosophie und -Qualität}
SignalReport verwendet eine gezielte Test-Strategie:
\begin{itemize}
    \item \textbf{Kernfunktionalität priorisiert}: Tests fokussieren auf kritische Geschäftslogik (Statistik-Berechnung, IP-Tracking, Measurer-Polymorphie), nicht auf Debug-Hilfen wie \texttt{toString()}
    \item \textbf{Realistische Szenarien}: Tests simulieren echte Nutzung (z. B. IP-Wechsel mit korrekter Zeitreihen-Sortierung, Wartungsfenster über Mitternacht)
    \item \textbf{Wartbarkeit}: Tests vermeiden fragile Assertions (z. B. exakte String-Formate) zugunsten stabiler, aussagekräftiger Prüfungen
    \item \textbf{100\% Bestehensquote}: Alle 156 automatisierten Tests laufen reproduzierbar grün -- nachweisbare Qualität für die Prüfung
\end{itemize}
